% titel: Blatt 10, A7: Wetterballon (Bonus, Ausblick Wärmelehre)
% skript: Kap. Ideales Gas, Dichte, barometrische H"ohenformel, Auftrieb

% Wetterballon: pV=nRT, Dichte rho=Mp/RT, barometrische H"ohenformel,
% Auftriebsgleichgewicht am Abheben und am Maximalvolumen.

Ein Wetterballon ist ein Paradebeispiel daf"ur, wie sich ideales Gas,
Dichte und Auftrieb zu einem praktischen Problem verbinden. Der Ballon
steigt, solange der Auftrieb sein Gewicht "ubersteigt; gleichzeitig f"allt
der Luftdruck mit der H"ohe, der Ballon bl"aht sich auf und wird dabei
d"unner.

\begin{werkzeugkasten}[Ideales Gas: Druck, Temperatur, Dichte]
Die allgemeine Gasgleichung
\begin{equation}
  p\,V=n_M R\,T
\end{equation}
(mit $n_M$ Stoffmenge, $R=\qty{8.314}{J/(mol\,K)}$). Aus der Stoffmenge
$n_M=m/M$ folgt die \emph{Massendichte}
\begin{equation}
  \rho=\frac{m}{V}=\frac{p\,M}{R\,T}.
\end{equation}
Jedes ideale Gas: dieselbe Temperatur/Druck $\Rightarrow$ \emph{gleiche
Zahl von Teilchen pro Volumen}, aber die Dichte skaliert mit der molaren
Masse $M$. Deshalb ist Helium ($M=\qty{4}{g/mol}$) viel leichter als Luft
(effektiv $\approx\qty{29}{g/mol}$) und tr"agt den Ballon.

\textbf{Effektive molare Masse der Luft} aus den Anteilen der
Hauptbestandteile:
\begin{equation}
  M_{\mathrm{Luft}}=0.78M_{\mathrm{N}_2}+0.21M_{\mathrm{O}_2}
  +0.01M_{\mathrm{Ar}}.
\end{equation}
\end{werkzeugkasten}

\begin{werkzeugkasten}[Auftrieb und das Loslösen des Ballons]
Ein eingetauchter Körper erf"ahrt einen Auftrieb gleich dem Gewicht der
verdr"angten Luft (Archimedes):
\begin{equation}
  F_{\text{Auftrieb}}=\rho_{\mathrm{Luft}}\,V_{\mathrm{Ballon}}\,g.
\end{equation}
Der Ballon beginnt zu steigen, sobald
\begin{equation}
  \rho_{\mathrm{Luft}}\,V\,g \ge (m_{\mathrm{Hülle}}+\rho_{\mathrm{He}}\,V)\,g,
\end{equation}
d. h. wenn die verdr"angte Luftmasse die Ballonmasse (Hülle + Helium)
gerade "ubersteigt. Physikalisch ist der Ballon ``leichter als Luft'',
weil $\rho_{\mathrm{He}}<\rho_{\mathrm{Luft}}$: das gef"ullte Helium wiegt
weniger als die Luft, die es verdr"angt. Um das Loslösen zu finden, l"ost
man nach dem n"otigen Volumen $V$ auf.
\end{werkzeugkasten}

\begin{werkzeugkasten}[Barometrische H"ohenformel]
Der Luftdruck nimmt mit der H"ohe $h$ exponentiell ab, weil die "ubereinander
liegende Luft ihr eigenes Gewicht tr"agt:
\begin{equation}
  p(h)=p_0\,e^{-M_{\mathrm{Luft}}\,g\,h/(R\,T_0)}.
\end{equation}
Die Skalenh"ohe $H=R T_0/(M_{\mathrm{Luft}} g)$ ist die H"ohe, in der der
Druck um $e^{-1}$ gefallen ist ($H\approx\qty{8}{km}$ f"ur die
Erdatmosph"are). Bei konstanter Temperatur und konstanter Heliummenge im
Ballon (die Hülle sei dicht, das Helium verhält sich wie ein ideales Gas)
bleibt $p\,V$ konstant, also
\begin{equation}
  p(h)\,V(h)=\mathrm{const}\quad\Rightarrow\quad
  V_{\mathrm{Ballon}}(h)=V_0\,\frac{p_0}{p(h)}=V_0\,e^{M g h/(R T_0)}.
\end{equation}
Je h"oher, desto kleiner der Druck, desto gr"o"ser das Volumen des
gleichbleibend gef"ullten Heliums. Der Ballon bl"aht damit exponentiell auf
und erreicht bei $V_{\max}$ sein Fassungsvermögen.
\end{werkzeugkasten}

\begin{achtung}[Konstanten und Grenzen]
- Normbedingungen: $T_0=\qty{273.15}{K}$, $p_0=\qty{1.013e5}{Pa}$
  (hier $T_0$ "uberall konstant angenommen).
- Molaren Massen in \emph{g/mol} umrechnen ($M\div1000$) vor dem Einsetzen
  in SI (kg/mol).
- Die Hülle sei sehr dünn: Eigengewicht $m_{\mathrm{Hülle}}=\qty{10}{kg}$;
  das Helium im Ballon hat \emph{denselben Druck} wie die Umgebung, so dass
  die Dichte in jedem H"ohepunkt aus dem dortigen $p$ folgt.
- Die barometrische Formel setzt konstante Temperatur voraus und gilt nur
  f"ur ein isothermes Modell der Atmosph"are --- in Wahrheit k"uhlt die
  Atmosph"are mit der H"ohe ab. Als Ausblick auf die Wärmelehre reicht das
  isotherme Modell f"ur die Gr"o"senordnung.
\end{achtung}
